% 論文節の日本語訳（読み物 PDF）用 preamble
% lualatex + jlreq
\documentclass[12pt,paper=a4]{jlreq}
\usepackage{luatexja-fontspec}
\usepackage{amsmath,amssymb,bm}
\usepackage{graphicx}
\usepackage{hyperref}
\usepackage{enumitem}
\usepackage{booktabs}
\usepackage{microtype}
\usepackage{xcolor}
\usepackage{caption}
\usepackage{tcolorbox}
\tcbuselibrary{breakable,skins}

\hypersetup{
  colorlinks=true,
  linkcolor=blue!60!black,
  urlcolor=blue!60!black,
  citecolor=blue!60!black,
}

\graphicspath{{images/}{./}}

\setlength{\parindent}{1em}
\setlength{\parskip}{0.55em}
\setlist[itemize]{leftmargin=1.6em,itemsep=0.25em,topsep=0.35em}
\setlist[enumerate]{leftmargin=1.8em,itemsep=0.25em,topsep=0.35em}

\captionsetup{
  font=small,
  labelfont=bf,
  skip=0.4em,
  justification=raggedright,
  singlelinecheck=false,
}

\newcommand{\vx}{\bm{x}}
\newcommand{\vu}{\bm{u}}
\newcommand{\veps}{\bm{\epsilon}}
\newcommand{\vw}{\bm{w}}

\DeclareMathOperator*{\argmax}{arg\,max}
\DeclareMathOperator*{\argmin}{arg\,min}
\DeclareMathOperator{\Tr}{Tr}
\DeclareMathOperator{\diag}{diag}

% 命題・定理ボックス（ページ跨ぎ時も切れ目を抑える）
\newtcolorbox{propositionbox}[1][]{%
  enhanced,
  breakable,
  colback=blue!4!white,
  colframe=blue!45!black,
  coltitle=white,
  fonttitle=\bfseries\small,
  title={#1},
  attach boxed title to top left={yshift=-2mm, xshift=3mm},
  boxed title style={
    colback=blue!45!black,
    colframe=blue!45!black,
    boxrule=0pt,
    arc=1pt,
  },
  boxrule=0.5pt,
  arc=1.5pt,
  left=5pt, right=5pt, top=8pt, bottom=5pt,
  before skip=0.9em,
  after skip=0.9em,
  pad at break=3mm,
}

% 図はフロートせず本文位置に置く（分断を防ぐ）
\newcommand{\translationfigure}[2]{%
  \begin{center}
    \includegraphics[width=0.88\linewidth]{#1}\\[0.35em]
    {\small #2\par}
  \end{center}
  \vspace{0.4em}%
}

\begin{document}
{\color{gray}\footnotesize 出典: §5.5\par}
\vspace{0.2em}
\begin{propositionbox}[原題]
\textbf{原題}: Closing Remarks
\textbf{日本語タイトル}: まとめ
\textbf{出典§}: §5.5 of *The Principles of Diffusion Models*
\end{propositionbox}


\section*{5.5 まとめ}


本章では、拡散モデルに対する第3の、そして最後の基礎的視点——決定論的フローの原理に根ざす視点——を明らかにした。探究は正規化フロー（Normalizing Flows, NFs）から始まり、変数変換公式を活用して単純な事前分布とデータ分布の間の厳密で可逆な写像を学習する枠組みを見た。続いてこの概念がニューラルODEへと連続時間過程として発展し、学習された速度場が変換を支配することを見た。しかしこの手法には、学習ループ内で高価なODEシミュレーションを要する、という重大な欠点がある。
現代のフローマッチング（Flow Matching, FM）の枠組みは、この課題に対する優雅で効率的な解として提示された。確率経路 $\{p_t\}_t$ と、連続の式を満たす対応する速度場をあらかじめ定義することで、FMはODEフローに対する明確な目標を確立する。決定的に重要なのは、変分的視点やスコアに基づく視点で見たのと同様に、FMが強力な条件付けの技法を用いることである。これにより、周辺速度場を整合させるという扱いづらい問題が、既知の条件付き速度に対する単純で扱いやすい回帰へと変換され、学習が完全にシミュレーション不要になる。この視点は、拡散モデル自体を、ガウス事前からデータ分布へ輸送する決定論的フローを学習する特別な場合として再解釈する。
フローに基づく視点の導入により、拡散モデリングの3つの概念的支柱に関する概観は完結した。この旅を通じて顕著なパターンが現れた：各枠組みは、VAE、EBM、あるいはNFという異なる起源をもちながらも、連続時間の生成過程に収束し、扱いやすい学習を可能にするために条件付け戦略に依拠してきた。
次章では、これらの並行する糸を、単一の統一的枠組みへと統合する。具体的には：

\begin{enumerate}
\item 変分的、スコアに基づく、フローに基づく視点が、単なる類似ではなく、基礎的なレベルで数学的に同値であることを形式的に示す。
\end{enumerate}


\begin{enumerate}
\item Fokker–Planck方程式が3つの視点すべてにわたる密度発展を支配する普遍法則として働き、それらが同じ核となる生成原理を記述する異なるレンズにすぎないことを明らかにする。
\end{enumerate}


この統一的レンズは、現代の拡散パラダイムに対する完全で体系的な理解を与える。

---

\section*{6 拡散モデルへの統一的・体系的レンズ}


\begin{propositionbox}[注意]
Mathematics is the art of giving the same name to different things.
— Henri Poincaré
\end{propositionbox}


本章は、変分的、スコアに基づく、フローに基づく視点を一貫した図像の中で接続する体系的観点を提示する。異なる直観に動機づけられながらも、これらのアプローチは現代拡散モデルの根底にある同じ核心機構に収束する：周辺分布の経路をたどる前方腐敗過程を定義し、この経路を逆方向にたどる生成ダイナミクスの構築に必要な局所量を学習する。第2章から第5章に基づき、共通のレシピを観察する：周辺の経路をたどる前方腐敗過程を定義し、この経路に沿って単純な事前をデータ分布へ輸送する時間変化ベクトル場を学習する。
すべての視点にわたる鍵となる成分は §6.1 で導入される条件付けの技法であり、扱いづらい周辺目的を扱いやすい条件付き目的へ変換し、安定で効率的な学習へ至る。
§6.2 では学習目的を体系的に解析し、その本質的構成要素を同定し、変分的、スコアに基づく、フローに基づく観点で損失関数がどう定式化されるかを明らかにする。
§6.3 では、$x_t=\alpha_t x_0+\sigma_t\epsilon$ の形の任意のアフィン前方ノイズ注入が、標準線形スケジュール $x_t=(1-t)x_0+t\epsilon$ へ同値に変換できることを示す。さらに、ノイズ予測、クリーンデータ予測、スコア予測、速度予測などの共通パラメータ化は、勾配のレベルで交換可能である。したがってノイズスケジューラとパラメータ化の選択は、いずれも同じモデリング原理に従う。
最後に §6.4 が議論をまとめ、支配法則——Fokker–Planck方程式——を同定する。変分スキーム（離散時間デノイジング）、スコアに基づく手法（SDE定式化）、あるいはフローに基づく手法（ODE定式化）のいずれとして見ても、それぞれは同じ密度発展に従う周辺をもつ生成器を構成する。したがってFokker–Planck方程式は3つの観点すべてが尊重する普遍制約として働き、違いはパラメータ化と学習目的にのみ生じる。

\textbf{本章の後、あなたは次ができるようになる…}

\begin{itemize}
\item 変分的、スコアに基づく、フローに基づく学習の背後にある共通の条件付け原理を認識する：扱いづらい周辺目標は、扱いやすい条件付き目標を通じて学習できる；
\item 4つの共通予測型——クリーンデータ、ノイズ、スコア、速度——を理解し、前方スケジュールが非退化のときそれらのオラクル予測間を変換する；
\item 拡散学習における主な設計選択——前方ノイズスケジュール、予測目標、時間サンプリング分布、時間重み付け——を解きほぐし、どの変更が数学的再パラメータ化であり、どれがなお実用的最適化に影響し得るかを理解する；
\item 「同値性」の異なる概念を区別する：同じオラクル情報の表現、重み付けや加法定数を除く同値な母集団目的、同じ周辺密度経路の追跡、実用的学習またはサンプリングアルゴリズムとしての同一振る舞い；
\item 異なるアフィン確率経路が時間再パラメータ化と空間的リスケーリングを通じてどう関係するか理解し、これらの変換が何を保ち何を保たないかを追跡する；
\item 連続の式とFokker–Planck方程式が、拡散モデルの変分的、決定論的ODE、確率的SDE記述を接続する密度レベルの原理をどう提供するかを理解する。
\end{itemize}


\textbf{読者ルート.} 本章は本書の概念的要石である。最初の通読では、条件付け原理、4つの予測型間の関係、3つの視点を接続する連続の式／Fokker–Planck視点に焦点を当てよ。中心的教訓はすべての変換公式を暗記することではなく、どの同値性の概念が用いられているか、その同値性が実際に何を保存するかを問うことである。重み付け、スケジュール、パラメータ化選択の詳細解析は、数学的に同値な定式化が実用ではなお異なる振る舞いをし得る理由を理解するために再訪できる。
\end{document}
